% Bốn mô hình của Chương 4 đặt cạnh nhau (bản 2D cho ảnh 32x32x3). Shape sau mỗi khối đọc từ
% compare.json/trace, tức kết quả thật của một lượt forward trong notebook 06, không gõ tay.
\begin{tikzpicture}[
    b/.style={khoi, minimum width=26mm, minimum height=5.5mm, font=\scriptsize},
    p/.style={khoiX, minimum width=26mm, minimum height=4.5mm, font=\scriptsize},
    h/.style={khoiG, minimum width=26mm, minimum height=5mm, font=\scriptsize},
    s/.style={font=\tiny, text=black!60, anchor=west},
    m/.style={-{Stealth[length=1.4mm]}, ptitblue},
    tieu/.style={font=\small\bfseries\sffamily, align=center},
    y=7.5mm]
  % ---------------- M0 BasicCNN: hàng i của trace basic_2d
  \begin{scope}
    \node[tieu, text=mZero] at (0, 1) {M0 BasicCNN};
    \foreach \y/\r/\sty/\txt in {0/0/p/Input, -1/1/b/{Conv 3×3 + ReLU}, -2/3/p/{MaxPool 2×2},
        -3/4/b/{Conv 3×3 + ReLU}, -4/6/p/{MaxPool 2×2}, -5/7/b/{Conv 3×3 + ReLU}, -6/9/p/{MaxPool 2×2},
        -7/10/p/Flatten, -8/11/h/{FC + ReLU}, -9/13/h/{FC $K$}} {
      \node[\sty] (n\r) at (0, \y) {\txt};
      \node[s] at (1.35, \y) {\SL{compare/trace/basic_2d/\r/1}};
    }
    \foreach \u/\v in {0/1,1/3,3/4,4/6,6/7,7/9,9/10,10/11,11/13} \draw[m] (n\u) -- (n\v);
  \end{scope}
  % ---------------- M1, M2, M3: cùng khung, hàng i của trace <model>_2d
  \foreach \col/\key/\ten/\mau/\khoi in {1/vgg/{M1 VGGNet}/mOne/{khối thường},
      2/resnet/{M2 ResNet}/mTwo/{khối residual}, 3/seresnet/{M3 SE-ResNet}/mThree/{khối residual + SE}} {
    \begin{scope}[shift={(\col*3.95, 0)}]
      \node[tieu, text=\mau] at (0, 1) {\ten};
      \tikzset{bk/.style={b, fill=\mau!18, draw=\mau}}
      \foreach \y/\r/\sty/\txt in {0/0/p/Input, -1.25/2/bk/{2 × \khoi}, -2.5/3/p/{MaxPool 2×2},
          -3.75/5/bk/{2 × \khoi}, -5/6/p/{MaxPool 2×2}, -6.25/8/bk/{2 × \khoi},
          -7.5/10/p/{Global Avg Pool}, -9/11/h/{FC $K$}} {
        \node[\sty] (n\r) at (0, \y) {\txt};
        \node[s] at (1.35, \y) {\SL{compare/trace/\key_2d/\r/1}};
      }
      \foreach \u/\v in {0/2,2/3,3/5,5/6,6/8,8/10,10/11} \draw[m] (n\u) -- (n\v);
    \end{scope}
  }
  % Chú giải nội dung một khối
  \node[font=\scriptsize, align=left, anchor=north west, draw=gray!50, rounded corners=2pt, inner sep=4pt]
    at (-1.6, -9.8) {%
    \textcolor{mOne}{\textbf{khối thường}}: $\mathrm{ReLU}\circ\mathrm{BN}\circ\mathrm{conv}\circ\mathrm{ReLU}\circ\mathrm{BN}\circ\mathrm{conv}$ \quad
    \\ \textcolor{mTwo}{\textbf{khối residual}}: $y = \mathrm{ReLU}(F(x) + x)$ \quad
    \textcolor{mThree}{\textbf{+ SE}}: $y = \mathrm{ReLU}(\mathrm{SE}(F(x)) + x)$};
\end{tikzpicture}
